\section*{Заключение}
\addcontentsline{toc}{section}{Заключение}

В работе была построена двухэтапная модель прогноза ощущения счастья. Главная
методологическая проблема состояла в том, что состояния отсутствуют для всех
unknown-наблюдений. Поэтому вместо некорректного сценария
\code{реальные состояния $\rightarrow$ счастье} использовалась схема
\code{причины $\rightarrow$ восстановленные состояния $\rightarrow$ счастье}.

После VIF-фильтрации из причин были удалены два мультиколлинеарных общественных
признака. По связи со счастьем и качеству восстановления были выбраны четыре
состояния: социальная поддержка, технологический прогресс, кредитный оптимизм и
продовольственная безопасность. Среднее качество их восстановления по OOF
составило $R^2 = 0{,}9148$.

Нелинейные преобразования были полезны для восстановления состояний, особенно
для индекса кредитного оптимизма. Как одиночные добавления они давали небольшой
прирост, но банк преобразований позволил улучшить самое сложное из выбранных
состояний. Графики по сообществам показали выраженную групповую структуру
данных: сообщества отличаются по среднему уровню счастья и по областям значений
общественных причин. При этом на выбранном rowwise-разбиении лучшая финальная
модель использовала сферу занятости, но не использовала \code{community one-hot}.

Итоговая логистическая регрессия на предсказанных состояниях, причинах после
VIF-фильтра и сфере занятости достигла
\textbf{micro precision $=0{,}6063$}. Требуемый порог $0{,}6$ был преодолен.
Полученный прогноз для 600 unknown-наблюдений сохранен в ноутбуке как результат
финального обучения.
